\medskip

\begin{minipage}{0.5\linewidth}
On considère un cube ABCDEFGH et l'espace est rapporté au repère orthonormal
$\left(\text{A}~;~\vect{\text{AB}},~\vect{\text{AD}},~\vect{\text{AE}}\right)$.

Pour tout réel $m$ appartenant à l'intervalle [0~;~ 1], on considère les points $K$ et $L$ de coordonnées:

\begin{center}$K(m~;~0~;~0)$\quad et \quad $L(1 - m~;~1~;~1)$.\end{center}
\end{minipage}\hfill
\begin{minipage}{0.45\linewidth}
\psset{unit=0.85cm}
\begin{pspicture}(6.6,6.6)
\pspolygon(0.4,1)(4.2,0.5)(4.2,4.8)(0.4,5.3)
\psline(4.2,0.5)(6.6,1.4)(6.6,5.7)(4.2,4.8)
\psline(6.6,5.7)(2.8,6.2)(0.4,5.3)
\psline(5,5.9)(0.4,5.3)(1.6,0.82)
\psline[linestyle=dashed](5,5.9)(6.6,1.4)(1.6,0.84)
\psline[linestyle=dashed](0.4,1)(2.8,1.9)(6.6,1.4)
\psline[linestyle=dashed](2.8,1.9)(2.8,6.2)
\uput[dl](0.4,1){A} \uput[d](4.2,0.5){B} \uput[dr](6.6,1.4){C} \uput[ul](2.8,1.9){D}
\uput[ul](0.4,5.3){E} \uput[dr](4.2,4.8){F} \uput[ur](6.6,5.7){G} \uput[u](2.8,6.2){H}
\uput[dl](1.6,0.84){K} \uput[u](5,5.9){L}
\end{pspicture}
\end{minipage}

\medskip

\begin{enumerate}
\item Donner les coordonnées des points E et C dans ce repère.
\item Dans cette question, $m = 0$. Ainsi, le point L(1~;~1~;~1) est confondu avec le point G, le point K(0~;~0~;~0) est confondu avec le point A et le plan (LEK) est donc le plan (GEA).
	\begin{enumerate}
		\item Justifier que le vecteur $\vect{\text{DB}}\begin{pmatrix}1\\-1\\0\end{pmatrix}$ est normal au plan (GEA).
		\item Déterminer une équation cartésienne du plan (GEA).
	\end{enumerate}
\end{enumerate}

On s'intéresse désormais à la nature de C$K$E$L$ en fonction du paramètre $m$.

\begin{enumerate}[resume]
\item Dans cette question, $m$ est un réel quelconque de l'intervalle [0~;~1].
	\begin{enumerate}
		\item Démontrer que C$K$E$L$ est un parallélogramme.
		\item Justifier que $\vect{K\text{C}} \cdot \vect{K\text{E}} = m(m - 1)$.
		\item Démontrer que C$K$E$L$ est un rectangle si, et seulement si, $m = 0$ ou $m = 1$.
	\end{enumerate}
\item Dans cette question, $m = \dfrac12$.Ainsi, L a pour coordonnées 
$\left(\dfrac12~;~1~;~1\right)$ et K a pour coordonnées $\left(\dfrac12~;~0~;~0\right)$.
	\begin{enumerate}
		\item Démontrer que le parallélogramme CKEL est alors un losange.
		\item À l'aide de la question \textbf{3.~b.}, déterminer une valeur approchée au degré près de la mesure de l'angle $\widehat{\text{CKE}}$.
	\end{enumerate}
\end{enumerate}
